\section{From Stridon to the Greek East}\label{sec:jerome-youth}

\subsection{A borderland family and a Roman education}

Stridon stood somewhere near the late Roman boundary of Dalmatia and Pannonia.
Jerome says Goths destroyed it (\emph{De viris illustribus} 135), a brief
notice that joins local loss to the wider instability of the fourth-century
Balkans. The exact site remains unknown. His family could finance advanced
education and later appears to have had property; Jerome's brother Paulinian
and a sister also entered Christian ascetic networks. The family was Christian,
but Jerome, following a common fourth-century pattern, was not baptized as an
infant.

At Rome he studied under the grammarian Aelius Donatus, mastering the canon of
Latin prose and poetry, rhetorical argument, memory, and style. He copied or
acquired books and visited places associated with the martyrs. His letters
later show a writer who could never simply cease being classically trained.
Even denunciations of Cicero employ Ciceronian resources. Education furnished
the precision and aggression with which he would defend asceticism and
Scripture.

The disputed birth date changes how this education is plotted. A birth in 331
makes Jerome considerably older at Rome, Antioch, and death; a birth around
347 makes his references to school contemporaries and recent imperial history
easier to coordinate. The older chronology traditionally appeals to Prosper,
whose exact entry still requires collation for this edition; the later date
should not be called unanimous. The biographical sequence is firmer than the
age labels.

Jerome was baptized at Rome, often placed during the episcopate of Liberius
(352--366). He later traveled to Trier, perhaps in connection with an
administrative career, and moved within the ascetic circle at Aquileia. There
he formed or deepened relationships with Rufinus, Bonosus, Heliodorus, and
others. Jerome later called the group a choir of the blessed, but conflict
eventually scattered it. Ascetic commitment grew within friendship and
correspondence, not in a social vacuum.

\subsection{Antioch, fever, and the Ciceronian dream}

Travel east brought Jerome to Antioch. Severe illness and the deaths of
companions marked the period. In Letter 22.30, written years later to the Roman
ascetic Eustochium, he tells of a feverish dream in which he stood before the
Judge and was accused of being “a follower of Cicero and not of Christ.” He
swore off pagan books, was beaten, and awoke bearing the memory of terror.

The dream is A, but its first surviving telling belongs to 384 and to a
deliberately rigorous instruction on virginity. It cannot be read as a diary
entry from the fever. Jerome's later writing continued to display classical
learning, but the reciprocal dossier needed to say how an opponent used the
dream was not inspected for this edition. The episode therefore reveals
ascetic conscience and the hazards of publishing the self as an example,
without carrying a claim about Rufinus's exact argument.

Jerome withdrew to the region of Chalcis east of Antioch around 375--377.
His letters do not describe a wholly isolated retreat. They disclose a
networked existence: books, copyists, visitors, ecclesiastical factions, and
difficult study. He began serious Syriac and Hebrew work with teachers,
describing the painful acquisition of unfamiliar sounds in Letter 125.12.
The degree of his later Hebrew command remains debated, but the labor was
neither invented nor wholly independent of Jewish and Greek scholarly help.

Chalcis also entangled him in the Antiochene schism. Rival claimants and
theological vocabularies divided the church. Letters 15 and 16 appeal to Pope
Damasus for guidance about communion and terminology. Jerome presents his
loyalty to Rome sharply, but the local conflict cannot be reduced to a clear
Roman answer confronting simple heresy. Around 378/379 Paulinus ordained him
a presbyter. Jerome seems to have accepted on conditions consistent with his
ascetic vocation and did not become an ordinary parish pastor.

\subsection{Constantinople and Greek Christian learning}

In Constantinople Jerome encountered Gregory of Nazianzus, whom he later
called his teacher. The city placed him close to the theological world of the
Council of Constantinople in 381, although claims about his exact official role
should remain modest. He read and translated Greek works, including Origen,
Eusebius, and homiletic materials. Origen's exegetical learning and spiritual
interpretation deeply marked him before the later controversy made that debt
dangerous.

This Greek apprenticeship complicates the later emblem of Jerome as an
isolated champion of Hebrew truth against a Greek Bible. His Latin biblical
work drew on a layered scholarly world: existing Old Latin forms, Greek
manuscripts, Origen's Hexaplar comparisons, Hebrew texts, Jewish teachers,
Christian predecessors, scribes, and patrons. His particular achievement was
real, but it was never creation from nothing.

\subsection{Ascetic vocation and social position}

Jerome's asceticism renounced marriage, sexual activity, status display, rich
food, and conventional career. It did not eliminate class, literacy, or
patronage. Books, long travel, buildings, teachers, and secretarial work
required resources. Elite Roman women would later finance and intellectually
sustain the Bethlehem enterprise. Jerome's insistence on renunciation could
therefore expose aristocratic wealth to radical demands while continuing to
depend on it.

His rhetoric often ranked virginity above marriage and used bodily discipline
as a battlefield of salvation. The hierarchy had patristic precedents, but
Jerome made it combative. The biographer must hold together three facts:
voluntary celibacy opened real alternatives to ordinary elite marriage;
women such as Marcella, Paula, and Eustochium exercised serious scriptural and
institutional agency within it; and polemical exaltation of virginity could
demean married Christians and turn the body into a theater of suspicion.

\evidenceaxis{Jerome lived ascetically in the region of Chalcis and pursued
language study there.}
{Letters 14--17 and the retrospective account of Hebrew study in Letter
125.12.}
{“Desert” names an ascetic geography, not total isolation. The letters reveal
networks, books, messengers, teachers, and ecclesiastical controversy.}
